\chapter{Descriptive Accuracy: Does Masking the Top Region Change the Answer? (Day 5)}
\label{chap:day05}

\section{Motivation}

Day 4 produced, for every sample, a ranked list of candidate explanation regions and a masking
primitive capable of blacking out the top-ranked one. Neither is useful on its own: ranking a region
first and being able to mask it says nothing yet about whether that region actually mattered to the
model's judgment. Descriptive Accuracy closes that gap. It is one of the six metrics defined by
\citet{arreche2024exai} for evaluating black-box explainable-AI frameworks, and the idea behind it is
simple and falsifiable: if the explanation method has correctly identified the region the model's
decision depends on, removing that region should be capable of changing the decision; if masking the
region never changes anything, the region was not load-bearing for the answer in the first place, and
the explanation pointing to it was not describing anything the model actually used. Day 5's task was
to run this test, for the first time in the pilot, against the redesigned proxy and the standardized
regions that Days 3 and 4 built specifically to support it.

\section{Implementation}

\texttt{descriptive\_accuracy.py}'s \texttt{evaluate()} function implements the test directly on top
of the two modules already built: \texttt{inference.py}'s \texttt{answer\_rule} and
\texttt{regions.py}'s \texttt{mask\_region}. For each sample, it masks the top-1 standardized region
(black-fill) on the original image, reruns \texttt{answer\_rule}, and compares the result to the
unmasked baseline answer already on record from Day 3. It then proceeds further along the same ranked
list: starting again from the original image, it masks the top-1 region \emph{and} the top-2 region
together, reruns \texttt{answer\_rule} a second time, and again compares to baseline. No new
pipeline stage was built for this chapter; Descriptive Accuracy is a composition of two modules that
already existed, run twice per sample with progressively more of the image blacked out.

\texttt{scripts/05\_descriptive\_accuracy.py} drives this evaluation over all 163 pilot samples,
reusing the \texttt{worker\_boxes} and \texttt{object\_boxes} already written to
\texttt{results/baseline\_predictions.csv} on Day 3 to reconstruct each sample's standardized regions
exactly as Day 4 ranked them, then calling \texttt{evaluate()} once per sample. Results are written to
\texttt{results/descriptive\_accuracy.csv}, one row per sample recording the baseline answer, the
masked answer under each masking level, and whether each masking level flipped the answer relative to
baseline. For the first ten samples whose top-1 masking flips the answer, the script additionally
renders an overlay image showing exactly which region was masked, saved to
\texttt{results/figures/descriptive\_accuracy/}, so that flip cases can be inspected visually rather
than taken only on the strength of a boolean column in a CSV.

\section{Results}

Across all 163 samples, masking the top-1 region flips the answer in \textbf{36.2\%} of cases.
Masking the top-1 and top-2 regions together flips the answer in \textbf{35.0\%} of cases --- a
\emph{lower} rate than masking top-1 alone, an anomaly examined in its own section below.

\begin{center}
\begin{tabular}{lr}
\toprule
\textbf{Masking level} & \textbf{Descriptive accuracy (answer flip rate)} \\
\midrule
Top-1 region & \textbf{36.2\%} \\
Top-1 + top-2 regions & 35.0\% \\
\bottomrule
\end{tabular}
\end{center}

\subsection{By primary class}

Flip rate under top-1 masking varies substantially by ground-truth \texttt{primary\_class}:
\texttt{struck\_by\_risk} flips most often, at 61.5\%, followed by \texttt{fall\_hazard} at 46.0\%;
\texttt{ppe\_violation} and \texttt{compliant} both flip at 28.0\%.

\begin{center}
\begin{tabular}{lr}
\toprule
\textbf{Primary class} & \textbf{Top-1 flip rate} \\
\midrule
struck\_by\_risk & 61.5\% \\
fall\_hazard & 46.0\% \\
ppe\_violation & 28.0\% \\
compliant & 28.0\% \\
\bottomrule
\end{tabular}
\end{center}

\subsection{By assigned rule}

Broken down by \texttt{assigned\_rule\_id} instead, top-1 flip rate is highest for rule~4 (44.0\%)
and rule~3 (42.9\%), lower for rule~2 (38.5\%), and lowest for rule~1 (27.0\%).

\begin{center}
\begin{tabular}{lr}
\toprule
\textbf{Assigned rule} & \textbf{Top-1 flip rate} \\
\midrule
rule\_4 (struck-by risk) & 44.0\% \\
rule\_3 (guardrail) & 42.9\% \\
rule\_2 (harness) & 38.5\% \\
rule\_1 (hard hat) & 27.0\% \\
\bottomrule
\end{tabular}
\end{center}

This per-rule ordering tracks Day 3's redesigned-proxy sensitivity ordering closely, and the
correspondence is not coincidental. Day 3 found rule~4's proximity check and rule~3's looser-threshold
guardrail check carrying the strongest real discriminative signal among the four rules, with rule~1's
hard-hat check the weakest (16.0\% sensitivity). The same ordering reappears here: a proxy that barely
discriminates between compliant and violating inputs in the first place also has comparatively little
for masking to disrupt, because so much of its output is already determined by something other than
the specific region being tested. A proxy with genuine discriminative signal, by contrast, is more
exposed to having that signal's load-bearing region removed, and so flips more often under masking.
Descriptive accuracy, in other words, is not an independent measurement here so much as a second view
of the same underlying proxy quality Day 3 already characterized.

\section{Worked Examples}
\label{sec:day05-worked}

Two of the ten saved flip overlays illustrate the range of behavior masking can produce, from the
mechanistically subtle to the intuitively clean.

\begin{figure}[h]
\centering
\includegraphics[width=0.85\linewidth]{descriptive_accuracy/0000069_rule_1_flip.png}
\caption{Image 69, rule~1 (hard hat): the top-1 region is the worker box, not the hard-hat box.
Masking it removes the worker from detection entirely, rerouting the answer through a fallback branch
rather than removing direct evidence about the hard hat --- the case traced in detail in
Section~\ref{sec:day05-anomaly}.}
\label{fig:day05-0000069}
\end{figure}

\begin{figure}[h]
\centering
\includegraphics[width=0.85\linewidth]{descriptive_accuracy/0000120_rule_4_flip.png}
\caption{Image 120, rule~4 (struck-by risk): the top-1 region is the excavator box itself. Masking it
removes the proximity check's other half, flipping \texttt{hazard} $\to$ \texttt{safe}.}
\label{fig:day05-0000120}
\end{figure}

Figure~\ref{fig:day05-0000120} is the intuitive case, and it is worth stating explicitly what makes it
intuitive: the masked region is exactly the rule-defining object, and removing it removes the hazard
signal in a way that matches what "explanation" should mean. Rule~4's proximity check depends on two
grounded objects, a worker and an excavator; masking the excavator box removes one of the two
quantities the proximity comparison needs, and the answer flips from \texttt{hazard} to \texttt{safe}
because the check can no longer establish that the worker is close to anything dangerous. This is
descriptive accuracy behaving exactly as the metric is designed to detect: the top-ranked region was,
in fact, the region the judgment depended on, and masking it changed the judgment for a reason that
tracks the underlying safety concept.

Figure~\ref{fig:day05-0000069} looks superficially similar --- an answer flip following a masked
top-1 region --- but the masked region is the \emph{worker} box, not the hard-hat box, and the
mechanism behind its flip is entirely different from Figure~\ref{fig:day05-0000120}'s. That mechanism
is the subject of the next section.

\section{The Non-Monotonicity Anomaly}
\label{sec:day05-anomaly}

Naively, masking \emph{more} of an image's evidence should never un-flip an answer that masking
\emph{less} had already flipped: if removing one region was enough to change the model's judgment,
removing that region and another should still leave the judgment changed, since nothing has been
added back. The headline numbers above show this expectation failing outright. Top-1+top-2 masking's
flip rate, 35.0\%, is \emph{lower} than top-1 masking's flip rate, 36.2\%, and the difference is not
sampling noise: in 15 of 163 samples, masking the second-ranked region in addition to the first
reverses a flip that masking the first region alone had already produced, restoring the original
baseline answer.

\subsection{The traced case: image 0000069, rule\_1}

Image \texttt{0000069} (rule~1, hard hat) was traced in full to establish the mechanism behind these
15 regressions concretely rather than dismissing them as noise.

\textbf{Baseline.} The sample's baseline detection returns two boxes: one worker box, with area
10{,}778~px\textsuperscript{2}, and one hard-hat box, with area 891~px\textsuperscript{2}. The worker
box is more than twelve times larger than the hard-hat box. Day 4's \texttt{standardize\_regions}
ranks candidate regions by area, descending, with no other signal available to rank by --- so the
worker box, not the hard-hat box, is ranked top-1, and the hard-hat box is ranked top-2.

\textbf{Masking top-1 (the worker).} With the worker box blacked out, Florence-2's worker-detection
call on the rerun returns no boxes at all: \texttt{worker\_boxes} comes back empty. This triggers the
\texttt{if not worker\_boxes} branch in \texttt{\_answer\_presence\_rule}
(\texttt{inference.py}), the same fallback path Chapter~3 introduced for samples with zero detected
workers: "no worker detected, so answer \texttt{compliant} if the safety object exists anywhere in the
frame, \texttt{violation} otherwise." The hard-hat box is still visible --- only the worker region was
masked --- so the fallback finds an object and answers \textbf{compliant}. Baseline's answer for this
sample was \texttt{violation}, so this is recorded as a flip.

\textbf{Also masking top-2 (the hard hat).} Starting over from the original image and masking both
the worker box and the hard-hat box together, neither region is visible to the rerun. The same
fallback branch fires again --- there is still no detected worker --- but this time the scene-level
existence check for the safety object also fails, because the hard hat is now hidden too. The fallback
therefore answers \textbf{violation}: the original baseline answer. Masking strictly more of the image
has reverted the answer to where it started.

\textbf{The mechanism, generalized.} This is a real interaction between the masking test and the
person-relative redesign's own fallback logic from Chapter~3, not a bug in either component
considered alone. Masking the \emph{worker} region does not simply remove one piece of evidence among
several that the per-worker coverage check weighs; it removes the precondition,
\texttt{worker\_boxes} being non-empty, that the per-worker check needs even to run. Once that
precondition fails, control silently reroutes to the cruder scene-level branch, whose answer then
depends only on whether the \emph{other} region (the safety object) happens to still be visible. That
dependency is binary and is not monotonic in how much of the image has been masked: visible-object
$\to$ \texttt{compliant}, masked-object $\to$ \texttt{violation}, regardless of whether one region or
two has been blacked out. Masking top-1 alone leaves the object visible and so reads as
\texttt{compliant}; masking top-1 and top-2 together hides the object too and so reads as
\texttt{violation} again --- and because the rule's binary answer space has only these two values,
hiding the object after the worker is already gone does not produce some third intermediate answer,
it produces the literal original baseline value. All fifteen regressions follow this exact
full-revert pattern: \texttt{masked\_answer\_top2} equal to \texttt{baseline\_answer}, not merely
unequal to \texttt{masked\_answer\_top1}\footnote{Recomputed directly from \texttt{results/descriptive\_accuracy.csv}: of the 15 samples with \texttt{answer\_changed\_top1=True} and \texttt{answer\_changed\_top2=False}, \texttt{masked\_answer\_top2} equals \texttt{baseline\_answer} for all 15, not 14 as an earlier informal note on this finding stated.}. Confirmed visually for \texttt{0000069} in
Figure~\ref{fig:day05-0000069}, the masked region drawn there is unambiguously the worker box, not the
hard-hat box, exactly as the traced mechanism requires.

\subsection{Why this is not a defect to patch}

It would be tempting to treat this as a bug to fix --- a "broken" fallback branch that should be
removed or special-cased. That would be the wrong response. The fallback branch exists because Day 3
needed some defined answer for the 12 of 163 samples with zero detected workers at baseline; it is not
incidental code, it is a deliberate, documented design decision serving a real and separate purpose.
The anomaly here is not that the fallback branch is wrong, but that the masking test, by construction,
can manufacture exactly the condition the fallback branch was built to handle --- zero detected
workers --- as a side effect of removing a region for an unrelated reason, namely testing whether that
region was load-bearing for the answer. The two pieces of logic are each correct on their own terms;
the anomaly is what happens when one test's side effect collides with the other's precondition.

\section{Standard vs. Non-Standard Choices}

\textbf{Standard}: deletion- or masking-based descriptive accuracy is a well-established
explanation-fidelity pattern in the explainability literature, and is adopted here directly as one of
the six metrics defined by \citet{arreche2024exai} for evaluating black-box XAI frameworks --- that
work is the direct methodological source for the metric as implemented in this chapter, rather than a
method invented for this pilot.

\textbf{Non-standard}: discovering and mechanistically tracing the non-monotonic "masking more
un-flips the answer" anomaly to a fallback-branch interaction, rather than treating the 15 affected
samples as unexplained noise or excluding them from the headline numbers, is this chapter's
contribution. The anomaly was run down to a specific, concrete case (\texttt{0000069}) and a specific
line of code (\texttt{if not worker\_boxes} in \texttt{\_answer\_presence\_rule}), and reported as a
real interaction between two correct pieces of logic rather than papered over as measurement error.

\section{Implications for the Paper}

The headline 36.2\%/35.0\% pair should be reported together, not as two independent numbers but as a
single finding: that the second number is \emph{lower} than the first is the more important fact about
this chapter's results than either number's absolute size. We recommend the paper state the
non-monotonicity finding as a generalizable diagnostic, applicable beyond this one pilot: \textbf{watch
for non-monotonic masking results as a sign of hidden fallback logic.} Any masking- or
deletion-based explanation-fidelity test, run against any pipeline with conditional branches, is
capable of producing this same pattern whenever the masking procedure can remove the precondition for
one branch and silently reroute execution to another --- the non-monotonicity is not specific to
Florence-2, to construction safety, or to this particular fallback, but to the general shape of
testing-by-deletion against branching logic.

This implication is flagged forward explicitly because it recurs later in the pilot under a different
name. Robustness (Chapter~9) and Bounded Completeness (Chapter~10) both rerun \texttt{answer\_rule} on
modified images and compare the result to baseline --- the same shape of test as this chapter, applied
to perturbations and to completeness checks rather than to masking. The same fallback-rerouting risk
applies there: a perturbation or mask that happens to remove the only detected worker, rather than the
safety object the rule is nominally about, can flip the answer for reasons disconnected from the safety
concept the rule is meant to test. We recommend that Chapters~9 and~10 check explicitly for this
pattern in their own results --- inspecting whether flips correlate with a worker box dropping out of
detection --- rather than reporting only an aggregate flip or completeness rate and treating any
non-monotonic or counter-intuitive cases within it as unexplained.
